\subsection{\large Archivo de Interfaz Gráfica \texttt{(interfaz.c)}}
Este archivo, de la misma manera que la conexión, se manipula en un proceso independiente (\textit{frontend}) para que no se presenten problemas de congelamiento o bloqueos inesperados en la vista mientras la red opera. Utiliza la biblioteca \texttt{GTK+} para el renderizado visual y emplea \texttt{Pthreads} para mantener la comunicación asíncrona con el \textit{backend}.

A continuación, se detalla el funcionamiento de los bloques principales y subrutinas que conforman la lógica de la interfaz que se muestra al usuario principal:

\begin{itemize}
    \item \textbf{\texttt{interfaz\_grafica()}:} Es la subrutina maestra del \textit{frontend}. Inicializa el entorno gráfico de GTK, carga los recursos visuales CSS, y levanta los dos hilos (\texttt{gui\_lector} y \texttt{gui\_escritor}). Posteriormente, estructura jerárquicamente los contenedores (\texttt{GtkBox}, \texttt{GtkScrolledWindow}) y widgets (\texttt{GtkEntry}, \texttt{GtkButton}), cediendo finalmente el control al bucle principal de eventos (\texttt{gtk\_main()}).
    
    \item \textbf{Hilos IPC (\texttt{hilo\_lector\_gui} e \texttt{hilo\_escritor\_gui}):} Operan en segundo plano sin interferir con el renderizado de la ventana. El hilo lector extrae los datos entrantes del Pipe A de forma bloqueante y delega la actualización visual al hilo principal de GTK mediante \texttt{g\_idle\_add}. El hilo escritor evalúa activamente la bandera volátil \texttt{gui\_escribe} para transferir los mensajes del usuario hacia el \textit{backend} a través del Pipe B.
    
    \item \textbf{Gestión Visual (\texttt{añadir\_burbuja} y \texttt{mostrar\_mensaje\_recibido}):} Se encargan de instanciar las etiquetas (\texttt{GtkLabel}) para el texto del chat. Estas funciones evalúan la procedencia del mensaje (\texttt{es\_mio}) para anclar la alineación a la derecha (nodo local) o a la izquierda (nodo remoto). Utilizan la sintaxis \textit{Pango Markup} para dar formato a las estampas de tiempo y vinculan los ID del CSS para colorear las burbujas.
    
    \item \textbf{Persistencia (\texttt{cargar\_historial}, \texttt{guardar\_en\_historial}, \texttt{limpiar\_historial\_pulsado}):} Implementan un mecanismo local de almacenamiento en el archivo de texto \texttt{historial\_chat.txt}. Permiten serializar (usando el formato \texttt{Origen|Milisegundos|Texto}) y deserializar las cadenas para restaurar conversaciones de sesiones previas al abrir la aplicación, así como truncar el archivo para vaciar el registro.
    
    \item \textbf{Gestores de Eventos (\textit{Callbacks}):} 
    \begin{itemize}
        \item \textbf{\texttt{boton\_pulsado()}:} Captura el texto del cuadro de entrada al presionar Enter o dar clic en "Enviar", renderiza la burbuja local, lo almacena en disco y levanta la bandera para su envío IPC.
        \item \textbf{\texttt{ventana\_destruida()}:} Intercepta el cierre de la ventana por parte del usuario, empaquetando un mensaje especial (\texttt{id\_emisor = -1}) para forzar el apagado en cascada del \textit{backend} y cerrando el ciclo de GTK.
        \item \textbf{\texttt{scroll\_al\_final()}:} Ajusta dinámicamente el límite superior del \texttt{GtkAdjustment} para que la vista descienda automáticamente al insertar una nueva burbuja de diálogo.
    \end{itemize}
    
    \item \textbf{\texttt{cargar\_css()}:} Función integradora que instancia un \texttt{GtkCssProvider} para interpretar la hoja de estilos \texttt{v-boton.css}, inyectando los gradientes, bordes suavizados y colores distintivos del diseño tipo Messenger directamente al \texttt{GdkScreen} de la aplicación.
\end{itemize}

\begin{lstlisting}[style=estiloC, caption={Archivo del Front-End (conexion P2P)}]
#include <gtk/gtk.h>
#include <string.h>
#include <stdbool.h>
#include <unistd.h>
#include <pthread.h>
#include "header.h"

// Variables de control exclusivas para regular el ciclo de vida de los hilos de la GUI
volatile int gui_activo = 1;
// Bandera volatil para notificar que el usuario presiono enviar y hay datos listos para el pipe B
volatile int gui_escribe = 0;

// Buffers globales estructurados para la transferencia asincrona de informacion en la GUI
Mensaje msg_para_enviar;
Mensaje msg_recibido_pipe;

// Parametros de configuracion de red por defecto para el nodo remoto (amigo)
int puerto_amigo = 4001;        // El puerto en el que el otro extremo escucha conexiones
char dir_amigo[] = "127.0.0.1"; // Direccion IP del nodo remoto (mapeada localmente por defecto)
int puerto_actual = 4000;       // Puerto asignado a nuestra instancia para desplegarlo en la cabecera

// Componentes globales de GTK requeridos para la manipulacion de la vista desde multiples funciones
GtkWidget *caja_mensajes;
GtkAdjustment *adj_chat;
char buffer[TAM_MAX];

// Puntero de referencia global para interactuar con los descriptores de la tuberia principal IPC
Tuberia *gui_pipe_ref;

// Funcion de callback de baja prioridad para ajustar la barra de desplazamiento de forma automatica
static gboolean scroll_al_final(gpointer data)
{
    // Extraemos el limite superior maximo del area de scroll dinamicamente
    double max = gtk_adjustment_get_upper(adj_chat);
    // Obtenemos las dimensiones proporcionales de la pagina visible actual
    double page = gtk_adjustment_get_page_size(adj_chat);
    // Movemos el scroll al fondo absoluto restando el tamaño de la pagina visible al maximo alcanzado
    gtk_adjustment_set_value(adj_chat, max - page);
    // Retornamos FALSE para indicarle al planificador de GLib que destruya este evento tras ejecutarse una vez
    return FALSE;
}

// Subrutina encargada de construir los elementos visuales (widgets) para representar los mensajes del chat
static void añadir_burbuja(const char *texto, gboolean es_mio, long long tiempo_ms)
{
    // 1. Instanciamos una etiqueta de GTK para alojar la cadena del mensaje enviado o recibido
    GtkWidget *burbuja = gtk_label_new(texto);
    // Habilitamos el ajuste de linea automatico si el texto supera las dimensiones horizontales maximas
    gtk_label_set_line_wrap(GTK_LABEL(burbuja), TRUE);
    // Configuramos que la ruptura de linea se haga por palabras y caracteres de forma limpia
    gtk_label_set_line_wrap_mode(GTK_LABEL(burbuja), PANGO_WRAP_WORD_CHAR);
    // Establecemos una restriccion de diseño fijando el ancho maximo en 35 caracteres por fila
    gtk_label_set_max_width_chars(GTK_LABEL(burbuja), 35);

    // 2. Declaramos y formateamos la cadena que mostrara la estampa de tiempo
    char str_tiempo[32];
    char markup_tiempo[128];
    formatear_tiempo_str(tiempo_ms, str_tiempo);

    // Formateamos usando sintaxis Pango Markup para reducir la fuente (font='8') y conservar un color oscuro
    sprintf(markup_tiempo, "<span font='8' foreground='#000000'>%s</span>", str_tiempo);
    GtkWidget *label_tiempo = gtk_label_new(NULL);
    // Inyectamos el formato enriquecido en la etiqueta correspondiente
    gtk_label_set_markup(GTK_LABEL(label_tiempo), markup_tiempo);

    // Alineamos verticalmente la etiqueta del tiempo al centro para mantener simetria visual
    gtk_widget_set_valign(label_tiempo, GTK_ALIGN_CENTER);

    // 3. Creamos un contenedor horizontal (HBox) para estructurar la linea del mensaje actual
    GtkWidget *alineacion = gtk_box_new(GTK_ORIENTATION_HORIZONTAL, 0);
    // Asignamos margenes externos a la burbuja para que no quede pegada a las paredes ni a otras burbujas
    gtk_widget_set_margin_start(burbuja, 15);
    gtk_widget_set_margin_end(burbuja, 15);
    gtk_widget_set_margin_top(burbuja, 5);
    gtk_widget_set_margin_bottom(burbuja, 5);

    // Condicional para ramificar el diseño estilo Messenger dependiendo de la procedencia del mensaje
    if (es_mio)
    {
        // Si es propio, alineamos todo el contenedor a la extrema derecha
        gtk_widget_set_halign(alineacion, GTK_ALIGN_END);
        // Le asignamos el ID CSS "mi-mensaje" definido en v-boton.css (azul con bordes suavizados)
        gtk_widget_set_name(burbuja, "mi-mensaje");

        // Estructuramos la distribucion: primero la estampa de tiempo, luego el texto a la derecha
        gtk_box_pack_start(GTK_BOX(alineacion), label_tiempo, FALSE, FALSE, 0);
        gtk_box_pack_start(GTK_BOX(alineacion), burbuja, FALSE, FALSE, 0);
    }
    else
    {
        // Si proviene de la red, alineamos el contenedor a la extrema izquierda
        gtk_widget_set_halign(alineacion, GTK_ALIGN_START);
        // Le asignamos el ID CSS "su-mensaje" definido en v-boton.css (rojo/guinda)
        gtk_widget_set_name(burbuja, "su-mensaje");

        // Estructuramos la distribucion: primero el texto, luego la estampa de tiempo a la derecha
        gtk_box_pack_start(GTK_BOX(alineacion), burbuja, FALSE, FALSE, 0);
        gtk_box_pack_start(GTK_BOX(alineacion), label_tiempo, FALSE, FALSE, 0);
    }

    // Insertamos la fila estructurada dentro de la caja contenedora vertical del chat
    gtk_container_add(GTK_CONTAINER(caja_mensajes), alineacion);
    // Forzamos el renderizado inmediato del nuevo componente y sus widgets secundarios
    gtk_widget_show_all(alineacion);

    // Agregamos la rutina de autoscroll al bucle principal de eventos para ejecutarse en el proximo ciclo de ocio
    g_idle_add(scroll_al_final, NULL);
}

// Funcion para persistir de forma local e incremental la bitacora de conversacion
void guardar_en_historial(const char *texto, gboolean es_mio, long long tiempo_ms) {
    // Abrimos el archivo en modo "a" (append) para añadir nuevas lineas al final del documento sin sobreescribir
    FILE *archivo = fopen("historial_chat.txt", "a"); 
    if (archivo != NULL) {
        // Almacenamos los datos serializados bajo el delimitador de tuberia: [Origen]|[Milisegundos]|[Texto]
        fprintf(archivo, "%d|%lld|%s\n", es_mio ? 1 : 0, tiempo_ms, texto);
        // Cerramos el flujo del archivo para asegurar la descarga fisica de datos en el disco
        fclose(archivo);
    }
}

// Callback intermedio invocado de forma segura por el hilo principal de GTK para procesar buffers remotos
static gboolean mostrar_mensaje_recibido(gpointer data)
{
    // Construimos graficamente la burbuja correspondiente a la izquierda utilizando los datos del pipe
    añadir_burbuja(msg_recibido_pipe.texto, FALSE, msg_recibido_pipe.tiempo);
    // Respaldamos de forma permanente el mensaje remoto en el archivo de texto local
    guardar_en_historial(msg_recibido_pipe.texto, FALSE, msg_recibido_pipe.tiempo);
    // Retornamos FALSE para liberar el callback del despachador principal de eventos
    return FALSE;
}

// Subrutina invocada al arranque para reconstruir la sesion anterior a partir del archivo de historial
void cargar_historial() {
    // Abrimos el archivo en modo lectura estricta
    FILE *archivo = fopen("historial_chat.txt", "r"); 
    // Si el archivo no existe en el directorio, abortamos la operacion limpiamente sin lanzar excepciones
    if (archivo == NULL) return; 

    char linea[TAM_MAX + 100]; 
    // Leemos secuencialmente linea por linea hasta llegar al final del archivo de texto
    while (fgets(linea, sizeof(linea), archivo) != NULL) {
        // Removemos de forma segura el caracter de salto de linea '\n' remanente de fgets
        linea[strcspn(linea, "\n")] = 0;

        int es_mio = 0;
        long long tiempo_ms = 0;
        char texto[TAM_MAX] = {0};

        // Encontramos el primer token de division '|' correspondiente al bit de procedencia
        char *ptr1 = strchr(linea, '|');
        if (ptr1) {
            *ptr1 = '\0'; // Dividimos la cadena temporalmente introduciendo un terminador nulo
            es_mio = atoi(linea); // Convertimos el fragmento ASCII a entero ordinario (0 o 1)
            
            // Encontramos el segundo token correspondiente al timestamp en milisegundos
            char *ptr2 = strchr(ptr1 + 1, '|');
            if (ptr2) {
                *ptr2 = '\0'; // Dividimos el fragmento final
                tiempo_ms = atoll(ptr1 + 1); // Convertimos la cadena de texto a entero largo de 64 bits (long long)
                strcpy(texto, ptr2 + 1);     // Extraemos el texto integro restante del mensaje
                
                // Pintamos la burbuja retrospectivamente preservando si era local o remota
                añadir_burbuja(texto, es_mio == 1 ? TRUE : FALSE, tiempo_ms);
            }
        }
    }
    // Cerramos el descriptor de lectura una vez concluida la carga inicial
    fclose(archivo);
}

// Callback de evento activado al pulsar el boton "Limpiar"
static void limpiar_historial_pulsado(GtkWidget *widget, gpointer data)
{
    // 1. Truncamos y vaciamos por completo el archivo de historial abriendolo en modo escritura "w"
    FILE *archivo = fopen("historial_chat.txt", "w");
    if (archivo != NULL) {
        fclose(archivo); // Al cerrarse de inmediato, el archivo queda vacio con 0 bytes
    }

    // 2. Iteramos sobre cada widget hijo contenido en la caja de mensajes y lo destruimos fisicamente
    gtk_container_foreach(GTK_CONTAINER(caja_mensajes), (GtkCallback)gtk_widget_destroy, NULL);
    
    // Desplegamos bitacora de control en consola
    printf("[GUI] Historial limpiado.\n");
}

// Callback invocado cuando el usuario hace clic en el boton Enviar o presiona Enter en la entrada de texto
static void boton_pulsado(GtkWidget *widget, gpointer data)
{
    GtkEntry *entry = GTK_ENTRY(data);
    // Recuperamos el buffer de texto plano de la caja de entrada de GTK
    const char *texto = gtk_entry_get_text(entry);

    // Proteccion elemental: solo procesamos la entrada si contiene caracteres validos
    if (strlen(texto) > 0)
    {
        // Registramos inmediatamente el tiempo del sistema en milisegundos para estampar la operacion
        long long mi_tiempo = get_time_ms();
        // Agregamos la representacion grafica del mensaje propio apuntando a la derecha
        añadir_burbuja(texto, TRUE, mi_tiempo);

        // Almacenamos localmente el mensaje en la bitacora persistente
        guardar_en_historial(texto, TRUE, mi_tiempo);

        // Empaquetamos y serializamos los datos crudos en la estructura Mensaje para su transmision IPC
        msg_para_enviar.id_emisor = getpid(); // Identificador unico del proceso emisor (Frontend)
        msg_para_enviar.tipo = ENVIO;         // Clasificamos la operacion como un mensaje de salida
        strncpy(msg_para_enviar.texto, texto, sizeof(msg_para_enviar.texto) - 1);
        strncpy(msg_para_enviar.ip_destino, dir_amigo, sizeof(msg_para_enviar.ip_destino) - 1);
        msg_para_enviar.puerto_destino = puerto_amigo;
        msg_para_enviar.tiempo = mi_tiempo;   // Guardamos el timestamp sincronizado para la transmision

        // Alzamos la bandera volatil para desbloquear la ejecucion del hilo escritor de la interfaz
        gui_escribe = 1;

        // Vaciamos la caja de texto para dejarla disponible para el proximo mensaje
        gtk_entry_set_text(entry, "");
    }
}

// Manejador del evento de destruccion de la ventana (Cierre de la aplicacion)
static void ventana_destruida(GtkWidget *widget, gpointer data)
{
    printf("[GUI] Cerrando ventana. Enviando señal de apagado...\n");

    // Preparamos un mensaje de control especial inyectando -1 en el emisor
    msg_para_enviar.id_emisor = -1;
    msg_para_enviar.tipo = SALIDA;
    strncpy(msg_para_enviar.texto, "Cierre de interfaz.", sizeof(msg_para_enviar.texto) - 1);
    
    // Forzamos una escritura bloqueante sincrona en el pipe B hacia el backend para ordenar su apagado masivo
    write(gui_pipe_ref->B[1], &msg_para_enviar, sizeof(Mensaje));
    
    // Bajamos las banderas globales de control para romper los ciclos infinitos de los hilos remanentes
    gui_escribe = 0; 
    gui_activo = 0; 
    
    // Rompemos el bucle principal de GTK para devolver la ejecucion al hilo principal
    gtk_main_quit();
}

// Cargador e interprete del archivo de diseño en cascada CSS para GTK
void cargar_css(void)
{
    // Instanciamos un proveedor de estilos CSS para GTK
    GtkCssProvider *provider = gtk_css_provider_new();
    // Recuperamos la pantalla por defecto asignada al entorno grafico actual
    GdkDisplay *display = gdk_display_get_default();
    GdkScreen *screen = gdk_display_get_default_screen(display);
    
    // Leemos e interpretamos los selectores declarados en "v-boton.css"
    gtk_css_provider_load_from_path(provider, "v-boton.css", NULL);
    // Aplicamos los estilos de forma global a la pantalla con prioridad de aplicacion para sobreescribir temas nativos
    gtk_css_provider_add_provider_for_screen(screen, GTK_STYLE_PROVIDER(provider), GTK_STYLE_PROVIDER_PRIORITY_APPLICATION);
    // Decrementamos el contador de referencias del proveedor para liberar memoria limpia
    g_object_unref(provider);
}

// === HILO 1 DE LA GUI: Lector dedicado en segundo plano para el Pipe A ===
void *hilo_lector_gui(void *arg)
{
    Tuberia *local = (Tuberia *)arg;
    printf("[GUI-HILO-LECTOR] Hilo de lectura de pipe iniciado...\n");

    // Se ejecutara de fondo mientras dure la sesion grafica sin interrumpir los hilos de render
    while (gui_activo)
    {
        Mensaje temporal;
        // Operacion bloqueante a nivel de Kernel: espera paquetes provenientes del backend en el pipe A[0]
        int bytes = read(local->A[0], &temporal, sizeof(Mensaje));

        // Validamos que se hayan extraido bytes coherentes y la interfaz siga en estado activo
        if (bytes > 0 && gui_activo)
        {
            // Si el mensaje entrante fue catalogado por el backend como una RECEPCION de red
            if (temporal.tipo == RECEPCION)
            {
                printf("Recibi mensaje de la conexion...\n");
                // Mapeamos los datos de forma segura en la estructura global dedicada a la GUI
                memcpy(&msg_recibido_pipe, &temporal, sizeof(Mensaje));
                // Delegamos de forma asincrona e indolora la insercion grafica en el hilo principal de GTK
                g_idle_add(mostrar_mensaje_recibido, NULL);
            }
        }
        else
        {
            // Si el descriptor es cerrado o da error, rompemos el ciclo de ejecucion de forma inmediata
            break;
        }
    }
    printf("[GUI-HILO-LECTOR] Saliendo del hilo lector de la interfaz.\n");
    // Terminamos de forma limpia e independiente la existencia de este hilo
    pthread_exit(NULL);
}

// === HILO 2 DE LA GUI: Escritor dedicado en segundo plano para canalizar en el Pipe B ===
void *hilo_escritor_gui(void *arg)
{
    Tuberia *local = (Tuberia *)arg;
    printf("[GUI-HILO-ESCRITOR] Hilo de escritura de pipe iniciado...\n");

    while (gui_activo)
    {
        // Evaluamos de forma constante mediante sondeo (polling) el estado de la bandera de salida
        if (gui_escribe == 1)
        {
            // Despachamos la estructura serializada por el pipe B para que el modulo de red lo procese
            write(local->B[1], &msg_para_enviar, sizeof(Mensaje));
            // Apagamos la señal inmediatamente para evitar reenvios duplicados
            gui_escribe = 0; 
        }
        // Dormimos el hilo durante 1 milisegundo para evitar que el bucle vacio consuma el 100% de un nucleo del CPU
        usleep(1000); 
    }
    printf("[GUI-HILO-ESCRITOR] Saliendo del hilo escritor de la interfaz.\n");
    pthread_exit(NULL);
}

// Inicializador maestro de la interfaz grafica de usuario
void interfaz_grafica(int argc, char *argv[], Tuberia *padre)
{
    // Descriptores internos para almacenar la configuracion geometrica y jerarquica de los componentes de GTK
    GtkWidget *window, *box, *entry, *scrolled_window, *hbox, *button, *header_box, *info_contacto;
    pthread_t gui_lector, gui_escritor;

    // Resguardamos localmente la direccion de memoria del canal de pipes cruzados
    gui_pipe_ref = padre;

    // Inicializamos el entorno basico de GTK abstrayendo los argumentos de linea de comandos
    gtk_init(&argc, &argv);
    // Inyectamos las hojas de estilo personalizadas estilo Messenger retro-futurista
    cargar_css();

    // --- CONCURRENCIA: LEVANTAR LOS DOS HILOS INDEPENDIENTES DE LA GUI ---
    pthread_create(&gui_lector, NULL, hilo_lector_gui, (void *)gui_pipe_ref);
    pthread_create(&gui_escritor, NULL, hilo_escritor_gui, (void *)gui_pipe_ref);

    // Construimos la ventana principal y definimos sus atributos generales
    window = gtk_window_new(GTK_WINDOW_TOPLEVEL);
    gtk_window_set_title(GTK_WINDOW(window), "MSN P2P - INTERFAZ");
    gtk_window_set_default_size(GTK_WINDOW(window), 450, 650);

    // Vinculamos la señal de destruccion fisica de la ventana al callback de control ventana_destruida
    g_signal_connect(window, "destroy", G_CALLBACK(ventana_destruida), NULL);

    // Contenedor base vertical que apilara la cabecera, la caja de texto y la seccion de entrada
    box = gtk_box_new(GTK_ORIENTATION_VERTICAL, 0);
    gtk_container_add(GTK_CONTAINER(window), box);

    // --- DISEÑO DE LA CABECERA (Header estilo MSN) ---
    header_box = gtk_box_new(GTK_ORIENTATION_HORIZONTAL, 10);
    gtk_widget_set_name(header_box, "header-messenger");
    gtk_box_pack_start(GTK_BOX(box), header_box, FALSE, FALSE, 0);
    info_contacto = gtk_label_new(NULL);

    // Redactamos el banner informativo inyectando etiquetas con formato enriquecido mediante snprintf
    snprintf(buffer, sizeof(buffer),
             "<span font='11' weight='bold' foreground='#003399'>MESSENGER P2P</span>\n"
             "<span font='9' foreground='#666'>Escucha: %d - Conectado</span>",
             puerto_actual);

    gtk_label_set_markup(GTK_LABEL(info_contacto), buffer);
    gtk_box_pack_start(GTK_BOX(header_box), info_contacto, FALSE, FALSE, 15);

    // --- COMPONENTE COMPLEMENTARIO: BOTON DE LIMPIEZA DE BITACORA ---
    GtkWidget *btn_limpiar = gtk_button_new_with_label("Limpiar");
    g_signal_connect(btn_limpiar, "clicked", G_CALLBACK(limpiar_historial_pulsado), NULL);
    // Empaquetamos el boton al final (pack_end) para arrastrarlo al extremo derecho de la cabecera
    gtk_box_pack_end(GTK_BOX(header_box), btn_limpiar, FALSE, FALSE, 15);

    // --- COMPONENTE DE DESPLAZAMIENTO (Scrolled Window) ---
    scrolled_window = gtk_scrolled_window_new(NULL, NULL);
    gtk_box_pack_start(GTK_BOX(box), scrolled_window, TRUE, TRUE, 0);
    // Extraemos la configuracion geometrica del ajuste vertical para controlar la posicion de la barra
    adj_chat = gtk_scrolled_window_get_vadjustment(GTK_SCROLLED_WINDOW(scrolled_window));

    // Instanciamos el contenedor interno vertical donde se iran apilando dinamicamente las lineas del chat
    caja_mensajes = gtk_box_new(GTK_ORIENTATION_VERTICAL, 0);
    gtk_widget_set_valign(caja_mensajes, GTK_ALIGN_END); // Anclamos el crecimiento hacia abajo
    gtk_container_add(GTK_CONTAINER(scrolled_window), caja_mensajes);

    // --- AREA INFERIOR DE CAPTURA (HBox de entrada y transmision) ---
    hbox = gtk_box_new(GTK_ORIENTATION_HORIZONTAL, 5);
    gtk_widget_set_margin_start(hbox, 10);
    gtk_widget_set_margin_end(hbox, 10);
    gtk_widget_set_margin_top(hbox, 10);
    gtk_widget_set_margin_bottom(hbox, 10);
    gtk_box_pack_start(GTK_BOX(box), hbox, FALSE, FALSE, 0);

    // Instanciamos el cuadro de entrada de texto plano (GtkEntry)
    entry = gtk_entry_new();
    gtk_entry_set_placeholder_text(GTK_ENTRY(entry), "Escribe un mensaje...");
    
    // Conectamos el evento "activate" (presionar Enter dentro del cuadro) al callback boton_pulsado
    g_signal_connect(entry, "activate", G_CALLBACK(boton_pulsado), entry);
    gtk_box_pack_start(GTK_BOX(hbox), entry, TRUE, TRUE, 0);

    // Instanciamos el boton fisico de Enviar
    button = gtk_button_new_with_label("Enviar");
    // Conectamos el click del raton al mismo callback compartiendo el puntero de la caja de texto
    g_signal_connect(button, "clicked", G_CALLBACK(boton_pulsado), entry);
    gtk_box_pack_start(GTK_BOX(hbox), button, FALSE, FALSE, 0);

    // Reconstruimos la conversacion anterior leyendo de disco antes de levantar la interfaz visible
    cargar_historial();

    // Mostramos la ventana y ordenamos la cascada jerarquica de widgets internos
    gtk_widget_show_all(window);
    
    // Bloqueamos el hilo principal cediendo el control total al despachador de eventos de GTK (Bucle principal)
    gtk_main();

    // Sincronizacion final: Al romperse el bucle, esperamos la terminacion e incorporacion de los hilos asincronos
    pthread_join(gui_lector, NULL);
    pthread_join(gui_escritor, NULL);
}
\end{lstlisting}